\section{P против NP: оплата сертификатов ранга}\label{sec:pnp}

Для близнецов (\cref{sec:reduction,sec:firstcause}) и для гипотезы Римана (\cref{sec:riemann})
зелёная машина довела нас до узла, который мы затем приняли декретом --- первопричиной \axiom{}.
Здесь, единственный раз среди ветвей, машина ведёт себя иначе: прочтение задачи $\mathrm{P}$ против
$\mathrm{NP}$~\cite{clay-pnp} на языке ранга есть \emph{зелёная теорема}, и декрет не нужен вовсе.

Прочтение физично. $\mathrm{NP}$-задача --- это \emph{полная оплата} всех сертификатов ранга:
решить поиск значит учесть каждого свидетеля по отдельности. $\mathrm{P}$ --- это \emph{быстрый
проезд} по рангу: движение двигателя, который пролетает ранг, не посещая всех его состояний.
Неравенство между ними в том, что быстрый проезд физически не может оплатить неограниченное
семейство сертификатов --- чистый пигеонхол. Всё это проверено машинно, без всякой аксиомы. Сразу
подчеркнём, чего мы \emph{не} утверждаем: это не классическое $\mathrm{P}\neq\mathrm{NP}$.
Абстрактный слой классов сложности пластичен (его можно скроить так, чтобы классы совпадали или
разделялись даром), поэтому ничего о настоящих классических классах из него не следует, и мы такой
заявки не делаем.

\subsection{Две стороны на языке ранга}

\begin{definition}[\code{RankFastTraversal} / \code{FullRankCertificatePayment}]\label{def:pnp-sides}
\emph{P-сторона} --- это \code{RankFastTraversal}: в ранжированном графе всякий легальный путь
\code{PathN} длины $n$ из $x$ удовлетворяет $n \le \operatorname{lexRank}(x)$. \emph{NP-сторона} ---
это \code{FullRankCertificatePayment}: существование инъективного конечноключевого сжатия
\emph{всей} семьи генеалогий, где каждый сертификат учтён индивидуально.
\code{UnboundedCertificateSupply} --- семья с бесконечным множеством индексов.
\end{definition}

P-сторону дарит сама конструкция, а не гипотеза.

\begin{theorem}[\code{rankFastTraversal_holds}]\label{thm:rankFastTraversal_holds}
\green{}\ Всякий ранжированный граф проезжает ранг быстро: для любого ранжированного графа $G$
всякий легальный путь \code{PathN} длины $n$ из $x$ удовлетворяет $n \le \operatorname{lexRank}(x)$.
\end{theorem}

\emph{Идея доказательства.} Прямое следствие того, что длина легального пути не превосходит
стартового ранга (\code{len_le_lexRank}, \cref{sec:engine}); проверка линейна по рангу. Проверить
один сертификат тоже всегда дёшево (\code{verificationEasy_always}, \green{}); дорого --- учесть их
всех сразу.

\subsection{Сердце неравенства: пигеонхол}

\begin{theorem}[\code{no_fullPayment_of_unboundedSupply}]\label{thm:no_fullPayment_of_unboundedSupply}
\green{}\ Быстрый конечноключевой двигатель не может полностью оплатить неограниченное семейство
сертификатов: если семья генеалогий $F$ удовлетворяет \code{UnboundedCertificateSupply}
(множество её индексов бесконечно), то \code{FullRankCertificatePayment} для $F$ не выполнено ---
инъективного конечноключевого сжатия $F$ не существует.
\end{theorem}

\emph{Идея доказательства.} Чистый пигеонхол. Конечный ключ, сжимающий бесконечную поставку, обязан
столкнуть двух разных свидетелей в один; иначе он различал бы бесконечно много объектов конечным
числом ярлыков. Столкновение --- это двойная бухгалтерия, невозможная оплата. Дёшево проехать ранг
и дорого коснуться каждого состояния, и второе несовместимо с первым при бесконечной поставке. Это
и есть зелёная теорема, из-за которой всё работает, и она ничем не обязана \axiom{}.

Осталось предъявить масштаб, где поставка действительно бесконечна, --- и он есть, без единой
гипотезы о близнецах.

\begin{theorem}[\code{concreteSupply_unbounded_smallScale}]\label{thm:concreteSupply_unbounded_smallScale}
\green{}\ При $A \le 4$ конкретная семья сертификатов бесконечна: тип индексов
\code{concreteFamily}\,$A\,1$ бесконечен, свидетель --- инъекция пятиадической цепи
\code{fiveAdicChainFlow} --- той самой цепи, что опровергла ветвь $A \le 4$ в редукции.
\end{theorem}

\begin{corollary}[\code{concrete_noFullPayment_smallScale}]\label{thm:concrete_noFullPayment_smallScale}
\green{}\ При $A \le 4$ полная оплата конкретной семьи невозможна:
\code{FullRankCertificatePayment}\,(\code{concreteFamily}\,$A\,1$) не выполнено.
\end{corollary}

\begin{proof}
Применяем \cref{thm:no_fullPayment_of_unboundedSupply} к бесконечной поставке
\cref{thm:concreteSupply_unbounded_smallScale}.
\end{proof}

\subsection{«NP $=$ полная оплата» как теорема}

Чтобы заголовок стал теоремой, а не образом, свяжем «полную оплату» с локальным P-успехом --- на
каждом масштабе.

\begin{theorem}[\code{concrete_localPSuccess_iff_fullPayment}]\label{thm:concrete_localPSuccess_iff_fullPayment}
\green{}\ Локальный P-успех эквивалентен полной оплате сертификатов: для любых $A, M_0$
\[
   \code{LocalPSuccess}\,(\code{concreteProblem}\,A\,M_0)
   \iff
   \code{FullRankCertificatePayment}\,(\code{concreteFamily}\,A\,M_0).
\]
\end{theorem}

\emph{Идея доказательства.} Обе альтернативы разрешения коллизии --- легальный цикл и невозможная
оплата --- уже сожжены зелёно (\code{no_extendedFlowResolutionAlternative}): резолвер не может
сжульничать, и единственный способ разрешить коллизии --- честно учесть каждого. Теперь неравенство
собирается в одно зелёное утверждение.

\begin{theorem}[\code{pnp_rank_separation_smallScale}]\label{thm:pnp_rank_separation_smallScale}
\green{}\ При $A \le 4$ одновременно выполнены пять утверждений:
\code{RankFastTraversal}\,(\code{concreteGraph}\,$A\,1$) (двигатель проезжает ранг быстро);
\code{VerificationEasy} для каждого сертификата (проверка легка);
\code{UnboundedCertificateSupply}\,(\code{concreteFamily}\,$A\,1$) (поставка неограничена);
$\lnot\,\code{FullRankCertificatePayment}\,(\code{concreteFamily}\,A\,1)$ (полная оплата
невозможна); и \code{LocalSearchIncompressible}\,(\code{concreteProblem}\,$A\,1$) (локальный поиск
несжимаем).
\end{theorem}

Пять граней --- быстрый проезд, лёгкая проверка, бесконечная поставка, невозможность полной оплаты,
несжимаемость --- сходятся в одном зелёном утверждении. Это «$\mathrm{P}\neq\mathrm{NP}$»
\emph{в ранговой модели}, доказанное без аксиом.

\paragraph{Раскол по масштабам.}
Это не противоречит декрету близнецов, который на своём масштабе как раз \emph{разрешает} коллизии.
Локальный P-успех и несжимаемость суть строгие отрицания, но живут на разных масштабах: декрет
работает при $A \ge 5$, где на своём масштабе первопричина полностью оплачивает сертификаты
(\code{decreedScale_fullPayment}, \yellow{}), а несжимаемость живёт при $A \le 4$, где поставка
бесконечна и оплата невозможна --- безусловно зелёно, ничем не обязано \axiom{}. Отметим асимметрию
с Риманом: там граница была \emph{нужна} для несущей леммы; здесь сепарация в ней не нуждается. Мы
оставляем асимметрию видимой, а не вставляем неиспользуемую гипотезу ради симметрии.

\subsection{Почему третьей границы декрета нет}

У близнецов и Римана декрет был честной второй границей. Естественно спросить о третьей,
$\mathrm{P}/\mathrm{NP}$-границе; ответ --- нет, и это доказано машинно. Все три мыслимые формы
такого поля --- зелёные вердикты: универсальная форма (во всяком добросовестном фрейме P-решатель
извлекает резолвер) \emph{опровержима} (\code{pnpLawUniversal_refuted}); decider-форма
(реконструкция резолвера из голого решателя) тоже \emph{опровержима}
(\code{pnpLawDeciderGated_refuted}), поскольку голый \code{PDecider} существует классически для
любого языка, а извлечённый из него резолвер бьётся о несжимаемость при $A \le 4$; а
экзистенциальная форма \emph{уже доказана} (\code{pnpLawExistential_green}), так что декретировать
её было бы пусто. Зеркало стоимости схлопывается само.

\begin{proposition}[\code{pnpLaw_asserts_separation}]\label{thm:pnpLaw_asserts_separation}
\green{}\ Закон $\mathrm{P}/\mathrm{NP}$ эквивалентен разделению \emph{без всякой границы}, и обе
стороны --- уже теоремы; значит, честного третьего поля декрета не существует.
\end{proposition}

\emph{Идея доказательства.} Абстрактный слой классов сложности пластичен: его можно скроить так,
что классы совпадают даром (\code{allPFrame}), и так, что разделяются даром (\code{constantsFrame}).
Попутно адверсариальный аудит вскрыл четвёртый эпизод вакуумности: голый \code{PDecider} не несёт
содержания сложности, поэтому над несжимаемым узлом decider-фронты
\code{FaithfulSelfReductionFront} и \code{CurrentExtractionFront} классически пусты, а их выводы о
разделении \warn{}~вакуумны. Затронут лишь decider-канал; \code{InP}-мост
\code{Step00ToClassicalBridge} остаётся честной условной формой. Подлинно недостающий объект --- не
пропозиция над абстрактными фреймами, а привязанная к данным реальная машинная модель --- тот же
урок, что дал спектральный аудит Римана.

\subsection{Эпистемическое дополнение}

Есть эпистемический двойник этого запрета, зеркалящий результаты о непознаваемости для Коллатца и
Янга--Миллса (\cref{sec:collatz}).

\begin{theorem}[\code{pnpCause_unknowable}]\label{thm:pnpCause_unknowable}
\green{}\ Внутреннее, на конечном топливе, решение $\mathrm{P}/\mathrm{NP}$ было бы вечным
двигателем: оно недостижимо изнутри системы, за тем же горизонтом, что и причина Коллатца.
\end{theorem}

\emph{Идея доказательства.} Зелёно, без декрета; таинт аксиомы не растёт. «Узнать нельзя изнутри»
здесь теорема, а не лозунг: изнутри модели с конечным топливом решение не самообосновать, ибо
попытка стоила бы ровно вечного двигателя (\cref{sec:engine}).

\paragraph{Честный статус.}
Среди ветвей это единственный случай, где неравенство --- зелёная теорема \emph{без} всякого
декрета: в ранговой модели быстрый двигатель доказуемо не оплачивает неограниченную поставку.
Термодинамическая интуиция --- полный учёт дороже быстрого движения --- доказана в ранговом мире;
её перенос в мир машин Тьюринга остаётся недостающим data-anchor. Классическое
$\mathrm{P}\neq\mathrm{NP}$~\cite{clay-pnp} не доказано и не объявляется.
